\begin{abstract}
RAG systems have become increasingly popular for enterprise applications as a way to reduce LLM hallucinations by grounding responses in retrieved documents. However, how these systems behave when the retrieved documents themselves are corrupted or semantically mutated has not been thoroughly studied. This work presents an empirical evaluation of RAG systems using 416 AI-generated semantic mutations. We use an LLM-as-a-Judge evaluation pipeline powered by Llama 4 Scout 17B, cross-validated with a human evaluator to ensure scoring reliability. By comparing the effects of model size (8B vs. 17B), the presence or absence of a defensive System Prompt, and public versus private data sources, we uncover several findings. First, larger models are actually more susceptible to misinformation when no System Prompt is provided. Second, RAG systems almost never hallucinate when tested on private data. And finally, across all runs, the hallucination rate never exceeded 0.24\%, effectively staying at zero.
\end{abstract}
